\documentclass[../main]{subfiles}
\begin{document}
\chapter{Electromagnetismo}
\img{images/01/electro.jpg}{Bobina de cobre}
\info{
Contenido}{
  Campo magnético creado por una corriente eléctrica, intensidad de
  campo, inducción,permeabilidad magnética,fuerza magnetomotriz,electroimanes.
}
\actividad{Resolver los siguientes problemas}
\begin{table}[ht]
  \caption{Relación entre intensidad de campo magnético $H$ e
  Inducción magnética $B$.}
  \begin{center}
    \begin{tabular}[c]{|p{2cm}|p{3cm}|p{3cm}|p{3cm}|}
      \hline
      & \multicolumn{3}{|c|}{$H(Av/m)$} \\
      \hline
      {$B(Teslas)$} & Hierro forjado & Chapa normal & Chapa Silicio \\
      \hline
      0,1 & 80 & 50 & 90 \\
      0,3 & 120 & 65&140 \\
      0,5&160&100&170\\
      0,7&230&180&240\\
      0,9&400&360&350\\
      1,1&650&675&530\\
      1,3&1000&1200&1300\\
      1,5&2400&2200&5000\\
      1,6&5800&3500&9000\\
      1,7&7000&5000&16500\\
      1,8&11000&10000&27500\\
      1,9&17000&16000&\\
      2&27000&32000&\\
      \hline
    \end{tabular}
  \end{center}
\end{table}
\begin{enumerate}
  \item ¿Cuál es la inducción magnética $B$ existente en la cara de
    un imán de sección circular del parlante de un equipo de audio
    con radio $r=3cm$ si es atravesado por un flujo de $\Phi = 8mWb$?
  \item ¿Cuál es el flujo magnético $\Phi$ que existe en el campo
    magnético producido por una bobina de una electroválvula si tiene
    esta un núcleo cilíndrico de diámetro $d=4cm$ y la inducción
    magnética en la misma es de $B=1,5T$?
  \item ¿Cuál sería la sección de un electroimán que teniendo una
    inducción de $B=1T$ produce un flujo magnético de $\Phi=10mWb$?
  \item Se desea fabricar un electroimán de fuerza magnetomotriz
    $\mathcal{F}=221Av$, si solamente se puede fabricar con un número
    entero de vueltas y corriente (ej.$1v,2v,3v \ldots nv$) (ej.$1A,
    2A, 3A\ldots nA$) ¿Cuáles son las $4$ únicas formas de fabricación posibles?
  \item Calcular la intensidad del campo $H$ en el interior de una
    bobina que tiene 400 vueltas, una corriente $I=9A$ y una longitud
    media de $L=60cm$.
  \item Si el campo magnético que se produce en una bobina de
    $L=80cm$ al circular $I=2A$ de corriente por ella es de
    $H=5000Av/m$, calcular el número de espiras que posee.
  \item Una intensidad de campo de $H=4000Av/m$ produce una f.m.m.
    igual a $500Av$. Averiguar la longitud media del núcleo.
  \item Calcular el campo magnético $H$ de un electroimán solenoide
    comercial de $24VDC$ que también funciona con $12VDC$. Si este en
    $24VDC$ funciona con $I=0.43A$ tiene una bobina de 1000 vueltas y
    38mm de longitud de núcleo¿Tendrá la misma intensidad de campo
    con ambas tensiones?
  \item Comercialmente se venden solenoides que son electroimanes de
    núcleo móvil, que tienen $5Kg$ de fuerza de agarre, y funcionan
    con $220V$. La fuerza que puede realizar está dada por la
    siguiente fórmula $ F = 40000 B^2 S$. Sabiendo que la sección es
    cuadrada de 2cm de lado. Determinar la inducción magnética $B$
    necesaria para que dicha fuerza sea posible.
  \item Hay electroimanes en lavadoras automáticas para accionar el
    cierre de válvulas de purga, estas requieren de fuerza de
    apertura como para cierre, el cierre es producido por un resorte
    antagonista conectado a un núcleo móvil que cuando se energiza
    vence la fuerza del resorte e impulsa al núcleo hacia dentro del
    solenoide solidario a la válvula. Si la electroválvula posee un
    núcleo de acero forjado de longitud de núcleo $L=2cm$ y sección
    $s=2cm^2$, una bobina de 200 vueltas y consume una corriente
    $I=0,58A$ ¿Cuál es la fuerza de accionamiento?
\end{enumerate}

Resumen de formulas

$ H=\frac{NI}{L}$; $ B=\mu H$; $\mathcal{F}=NI$;
$\mu_r=\mu/\mu_o$; $\mu_o= 4\pi10^{-7}H/m$; $ F = 40000 B^2 S$

\end{document}
